The preceding theorem describes what happens as the mass varies. We also need to compare two positions on the same board. A Lipschitz
bound controls the change in the profile in terms of the distance between the positions. The question is whether the same constant
works as fewer and fewer rows remain empty.

\begin{teorema}[Spatial regularity]\label{thm:lipschitz_intro} \begin{enumerate}[(a)] \item For $\rho\in(0,1)$ and
$\lambda=\lambda(\rho)$, the vacancies and projections are Lipschitz on their closed intervals with constant
$L(\lambda)=2\lambda+8\lambda^2(1+\lambda)$, and $\nu_\rho$ is Lipschitz on the square with constant $6\lambda L(\lambda)$ in the
maximum norm. The bounds are uniform on each compact subset of $(0,1)$. \item There exists $L_*<\infty$ such that, for every
$x\in[0,1/2]$, the rescaled factors $r_R,r_C,u_D,u_A$ are Lipschitz with constant $L_*$, and $\nu$ is bounded and Lipschitz with
constants independent of $x$. \end{enumerate} \end{teorema}

The bound in part (a) degenerates as $\rho\to1$ because $\lambda(\rho)\to\infty$. Part (b) lets us work all the way to the full board.

For $x\in(0,1)$ and $\rho=1-x$, define \begin{equation}\label{eq:F}
 F(x)=\frac{1-x}{\lambda(1-x)\,x^2}=\frac{1-x}{b}=\int_{[0,1]^2}r_R(r)\,r_C(c)\,u_D(r-c)\,u_A(r+c)\,dr\,dc .
\end{equation} The last expression is the integral of $u_D(r-c)\,u_A(r+c)$ with respect to the weight $r_R(r)\,r_C(c)$, a probability
density on the square. The significance of $F$ lies in the integral identity in the following theorem, which recovers $\sopt(1)$ from
it.

\begin{teorema}[The function $F$ and the optimal value]\label{thm:F_intro} The function $F$ extends to a positive real-analytic
function on a neighborhood of $[0,1]$, with \[
 F(0)=\frac1{b_0},\qquad F(1)=1,\qquad \int_0^1\log F(x)\,dx=1+\sopt(1).
\] Moreover, as $x\downarrow0$, \[
 \sopt(1-x)-\sopt(1)=2x\log\frac1x+O(x).
\] \end{teorema}

\begin{observacion}[Numerical values] A deterministic numerical solution of the Euler system (Section~\ref{sec:discusion}) gives
$F(0)\approx0.0809$, $F'(0)\approx0.516$, and $\int_0^1\log F\approx-0.9440$, so that $\sopt(1)\approx-1.9440$. This value is
consistent with known values of the constant $\alpha$~\cite{Simkin,NAB}, as one would expect if $\alpha=-\sopt(1)$ (see
Remark~\ref{obs:simkin}). These numerical values carry no certified error bounds. \end{observacion}

\subsection{From the maximizer to the limiting profile}

At mass one, existence and uniqueness of the entropy maximizer are given in Lemma~2.22 of Simkin~\cite{Simkin}. Remark~\ref{obs:simkin}
specifies the correspondence between the functionals. The family of partial masses lets us study how this maximizer is approached. We
first derive the Euler equation in the interior (Section~\ref{sec:euler}). Its rescaled bounds allow us to identify the factors that
persist at mass one (Sections~\ref{sec:cotas} and~\ref{sec:limite}). Section~\ref{sec:analiticidad} completes the analysis of the
dependence on mass. To control spatial variation, we compare the profile with its translates (Section~\ref{sec:regularidad}). Finally,
integrating the Euler equation gives the relation between $F$ and the optimal entropy (Section~\ref{sec:F}), which we evaluate
numerically in Section~\ref{sec:discusion}.

\subsection{Relation to previous work}

Simkin~\cite{Simkin} bounds the optimal value at mass one by discretizing the problem. His upper bound uses Lagrange duality, which
expresses the discrete maximizer as an exponential in the multipliers~\cite[Section~6]{Simkin}. Nobel, Agrawal, and Boyd~\cite{NAB}
solve a discretization of the same problem by Newton's method and obtain the interval quoted above. It is a classical principle that an
entropy maximizer under linear constraints has exponential form~\cite{Csiszar}. What this paper adds is the study of the family of partial
masses, with its Euler equation, and above all uniform control up to the full board: explicit bounds, convergence and analyticity
of the rescaled profile, uniform Lipschitz regularity, and the integral identity for $F$. To the best of our knowledge, these results are new. We make no claim of novelty for the general techniques of convex analysis, duality,
or the implicit function theorem~\cite{Brezis,Zeidler}.
